\scanpage{21}
\mspage{5)}
où
\[
(5)\quad Q_X=P_{X,x}\wedge_{\mathcal M_{g,\nu}}\mathcal N_{g,\nu},\qquad
R=P_{X,x}\wedge_{\mathcal M_{g,\nu}}\Gamma_{\overline{\mathbb Q}/\mathbb Q}
\simeq Q_X\wedge_{\mathcal N_{g,\nu}}\Gamma_{\overline{\mathbb Q}/\mathbb Q}
\]
Ainsi, $\widetilde{\mathbb Q}$ et (1) sont connus par la connaissance du
\inserted{$\mathcal M_{g,\nu}$-tors. $P_{X,x}$} (3). Mais bien sûr même si
$k$ est alg. sur $\mathbb Q$ ($k=\widetilde{\mathbb Q}$, voire $k=\overline{\mathbb Q}$
i.e. \unsure{avec} $R$ trivialisé) la donnée de \deleted{\unsure{Q}} $P_{X,x}$
ne permet pas de récupérer $X$ sur $k$ --- p.ex. les autom. du torseur $P_{X,x}$
sont les élémts de $\mathcal M_{X,x}$ --- ils ont la puissance du continu ---
ceux du \deleted{\unsure{Q}} $\widehat{\mathcal G}_{g,\nu}$-torseur correspondant
à une restriction du groupe structural à $\widehat\Gamma_{g,\nu}$ i.e. à
un iso. $\widetilde{\mathbb Q}\simeq\overline{\mathbb Q}$, correspondent aux
éléments de $\widehat{\mathcal G}_{X,x}$ --- il a la puissance du continu,
tandis que les autom. de $X,x$ compatibles avec l'identité sur
$k=\widetilde{\mathbb Q}$ ($\overset{\mathrm{donné}}\simeq\overline{\mathbb Q}$)
font un groupe \uline{fini} !

Mais\ldots
\[
(6)\qquad k=\underset{i}{\uline{\lim}}\ K_i
\]
les $K_i$ sous-corps de t.f. --- et \deleted{\unsure{une}}\inserted{la}
structure etc $(X,x)$ implique une structure
\mspage{20}
